\uuid{Zhzc}
\exo7id{4221}
\titre{exo7 4221}
\auteur{quercia}
\organisation{exo7}
\datecreate{2010-03-12}
\isIndication{false}
\isCorrection{false}
\chapitre{Calcul d'intégrales}
\sousChapitre{Théorie}
\cpge{mpsi-mp2i.limites-continuite.c, pcsi.limites-continuite.c, ptsi.limites-continuite.c, mpsi-mp2i.integration.c, pcsi.integration.b, ptsi.integration.b, bcpst1.limites-continuite.c, bcpst1.integration-segment.a, bcpst1.integration-segment.b}
\module{Analyse}
\niveau{L1}
\difficulte{}

\contenu{
\question{
Soit $f : {[a,b]} \to \R$ continue non identiquement nulle, telle que :
$\forall\ k \in \{ 0, 1, \dots, n-1 \},\  \int_{t=a}^b t^kf(t)\,d t = 0$.

Démontrer, par récurrence, que $f$ change au moins $n$ fois de signe
sur $]a,b[$ (raisonner par l'absurde).
}
}
